\documentclass[11pt,a4paper]{article}
\usepackage[T1]{fontenc}
\usepackage[utf8]{inputenc}
\usepackage{lmodern,microtype}
\usepackage[margin=25mm,headheight=15pt]{geometry}
\usepackage{amsmath,amssymb,amsthm,mathtools,bm}
\usepackage{booktabs,array,enumitem,needspace,tabularx,threeparttable}
\usepackage{graphicx,xcolor}
\usepackage{csquotes}
\usepackage[backend=biber,style=numeric,sorting=none,maxbibnames=20,giveninits=true,doi=true,url=false,isbn=false]{biblatex}
\addbibresource{../../theory/references.bib}
\usepackage[colorlinks=true,linkcolor=blue!45!black,citecolor=blue!45!black,urlcolor=blue!55!black]{hyperref}
\usepackage[nameinlink,noabbrev]{cleveref}
\usepackage{fancyhdr}
\pagestyle{fancy}
\fancyhf{}
\fancyhead[L]{\small Acoustic Freeform Lab}
\fancyhead[R]{\small S09 results report}
\fancyfoot[C]{\thepage}
\setlength{\parskip}{0.35em}
\setlength{\parindent}{1.2em}
\clubpenalty=10000
\widowpenalty=10000
\displaywidowpenalty=10000
\AtBeginBibliography{\small}
\setlist{itemsep=0.2em,topsep=0.3em}
\numberwithin{equation}{section}
\newcommand{\dd}{\,\mathrm d}
\newcommand{\um}{\,\textmu m}
\newcommand{\nm}{\,nm}
\def\UrlBreaks{\do\/\do-\do\_\do.\do,}
\renewcommand{\UrlFont}{\ttfamily\small}
\newcommand{\file}[1]{\nolinkurl{#1}}
\newcommand{\figdir}{../../../artifacts/studies/S09-universal-source}
\newcommand{\genfig}{../../../artifacts/reports/s09-results/figures}
\newcolumntype{R}{>{\raggedleft\arraybackslash}X}
\newcolumntype{L}{>{\raggedright\arraybackslash}X}
\DeclareMathOperator{\tr}{tr}
\DeclareMathOperator{\sym}{sym}
\DeclareMathOperator{\Rea}{Re}
\hypersetup{pdftitle={One machine, many stigmatic singlets: numerical results of study S09},pdfauthor={Acoustic Freeform Lab},pdfsubject={Acoustic shaping of a UV-curable liquid lens: results report}}
\title{One machine, many stigmatic singlets\\[0.5em]
\large Numerical results of study S09 on acoustic shaping of NOA~61 between two water layers}
\author{Acoustic Freeform Lab\\\small Results report for colleagues; private draft}
\date{27 October 2026}
\begin{document}
\maketitle
\begin{abstract}
We report the numerical results of study S09: whether one fixed acoustic machine
(a 4\,mm water / NOA~61 / water cell with fixed windows and 112 source channels)
can shape the two liquid interfaces into the faces of the stigmatic singlets of
Silva-Lora and Torres~\cite{silva2020}, with the lens selected only by the drive
command, the heights of two pinning edges and the fill volumes. Within a declared
axisymmetric, lossless-acoustics, quiescent creeping-flow model the answer is yes
for every one of the 68 family lenses that fit the cell: a convex certificate
bounds the reachable pupil error below 0.5\,nm, and 67 of the 68 are formed from
the liquid at rest within 1\,nm by one or two tones. Three designs passed the full
chain of independent numerical checks. All results are numerical; no machine has
been built and nothing has been measured. Sensitivities of tens to hundreds of
nanometres to a 0.1\,\% sound-speed change, and unmodelled acoustic streaming,
heating and three-dimensional disturbances, separate these numbers from a
physical lens.
\end{abstract}
{\small\tableofcontents}
\clearpage

%======================================================================
\section{Summary}\label{sec:summary}

\paragraph{Question.} With the cell, liquids and source layout fixed, which
stigmatic singlets (SOSLs) of the Silva-Lora and Torres family~\cite{silva2020}
can be formed faithfully from the liquid at rest, if the lens is chosen only by
the drive command, the two pin heights and the three fill volumes? Which extra
source freedom extends the set most cheaply?

\paragraph{Answer within the model.} Everything below holds only in the declared
numerical model of \cref{sec:model}, which is axisymmetric and has lossless longitudinal acoustics
with ideal velocity sources, time-averaged radiation stress as the only acoustic
load, quiescent liquids and creeping flow.
\begin{itemize}
\item \textbf{Capability.} Of the 81 family singlets scanned (scaled 1/14 to a
4\,mm aperture), 68 fit the common cell. For all 68 the convex capability
certificate is below 0.5\,nm within 1\,W (largest 0.47\,nm, lens 57 with one
band), so every lens that fits is reachable in principle by this source.
\item \textbf{Formation from rest.} Started from the acoustics-off meniscus, the
nonlinear start-up branch ends within 1\,nm of both faces for 61 lenses with the
first command class tried (51 with one tone at 1.5\,MHz, 10 with one tone at
0.75\,MHz plus one at 1.5\,MHz), and for 6 more once a second tone at 0.75\,MHz
is added: 67 of 68. Lens 57 is the one exception, and it is a failure of the
start-up path, not of capability.
\item \textbf{Full chain.} Fig.~9b, Fig.~7a with a static fold margin and family
lens 7 (two tones, f/0.79) passed fixed-command refinement, coupled linear
stability and formation in time at two steps, and form within 1\,nm
(\cref{tab:published}). For two rescued lenses (15 and 71) refinement and
stability are done; their formation runs are pending.
\item \textbf{Hardware.} A first-order channel study suggests that 14 channels
already keep the certificate at 0.025\,nm for five representative lenses.
\item \textbf{Design tool.} \texttt{afl design} reproduced the Fig.~9b result end
to end from the lens prescription alone (status \enquote{formed},
0.044\,/\,0.033\,nm).
\end{itemize}

\paragraph{Main numbers.} Fig.~9b on the common machine: one tone of 0.17\,W,
sampled peak pressure 0.55\,MPa, formed pupil error 0.05\,/\,0.04\,nm
(front\,/\,back), from an acoustics-off meniscus 0.77\,/\,2.07\um\ away; Fig.~7a
with fold margin 0.7: 1.12\,W, 2.10\,MPa, formed 0.53\,/\,0.41\,nm; lens 7: two
tones, 0.55\,W, formed 0.12\,/\,0.61\,nm from 1.10\,/\,43.8\um.

\paragraph{Main limits.} (i) A 0.1\,\% change of the water sound speed moves the
faces by 28--299\,nm (front) depending on the design (\cref{tab:sensitivity}); any
physical machine would need wavefront measurement and closed-loop trim, and the
NOA~61 sound speed is a placeholder. (ii) Acoustic streaming, heating,
non-axisymmetric disturbances, inertia, contact-line hysteresis, transducers and
curing are not modelled. (iii) The full chain has been run on three designs, the
start-up branch alone on all 68. (iv) The Fig.~7a fold-margin design uses 1.12\,W
and 2.10\,MPa, slightly above the machine's nominal 1\,W\,/\,2\,MPa caps.
(v) Nothing has been built or measured: this is numerical evidence for a model,
not experimental validation.

%======================================================================
\section{The model in brief}\label{sec:model}

The formulation is the project's theory document (LaTeX source
\file{docs/theory/}); the steps used here are its sections \enquote{Faithful
two-face shaping: reachable equilibria of a modal source}
(\file{sections/10e_faithful_inverse.tex}) and \enquote{A common source: load
demand and convex capability} (\file{sections/10f_universal_source.tex}). We do
not re-derive them. Acoustic radiation pressure deforming a liquid interface is
established experimentally and numerically~\cite{bertin2012,chesneau2022,sisombat2023};
fluidic shaping of optical components by other means is reported
in~\cite{elgarisi2021}. Neither source tests the multi-interface, multi-lens
machine considered here.

\paragraph{Cell.} A cylinder of pinning radius 2.02\,mm (clear aperture 4\,mm)
with windows at $z=-0.8$ and $+2.2286$\,mm, the front lens vertex at $z=0$, holds
three layers: water, uncured Norland NOA~61, water. Both interfaces are pinned on
edges at the lens rim heights; each layer is filled to the volume that the target
faces enclose (\file{configs/machine/v1.json}). The windows are the same for every
lens; only the pin heights and fill volumes follow the lens.

\paragraph{Liquids.} \Cref{tab:materials} lists the material record
\file{configs/shared/noa61-water-pressure-reference.json}. Its own status notes
say that the resin values are manufacturer typical values, the resin--water
interfacial tension is an engineering assumption, and the resin sound speed is a
placeholder.

\begin{table}[htbp]\centering\small
\caption{Material constants of the model (SI). Status as declared in the material record.}\label{tab:materials}
\begin{tabular}{@{}lrrr@{}}
\toprule
 & Water (lower, upper) & NOA 61, uncured & Status \\
\midrule
Density (kg\,m$^{-3}$) & 997 & 1231 & typical values \\
Sound speed (m\,s$^{-1}$) & 1497 & 1500 & resin: placeholder \\
Viscosity (Pa\,s) & 0.00089 & 0.30 & typical values \\
Interfacial tension, both faces (N\,m$^{-1}$) & \multicolumn{2}{c}{0.025} & assumption \\
Gravity (m\,s$^{-2}$) & \multicolumn{2}{c}{9.81} & \\
\bottomrule
\end{tabular}
\end{table}

\paragraph{Sources.} Ideal normal-velocity sources on both windows and the wall:
32 ring regions per window and 16 wall bands per layer (112 channels), driven
through the same 26 continuous low-order modes for every lens (8 Bessel modes
$J_0(j'_{0,k}r/R)$ per window and 10 cosines on the wall). The region count is a
quadrature of the continuous source space, not an apparatus property; the
channel study (\cref{sec:channels}) tests it. Drive options: one tone at
1.5\,MHz, or one tone each at 0.75 and 1.5\,MHz. Limits: 1\,W source power and
2\,MPa sampled peak pressure per band.

\paragraph{Physics implemented.} Time-harmonic, axisymmetric, lossless
longitudinal acoustics in the three layers with passive ports; the time-averaged
radiation stress on the two interfaces is the only acoustic load. Tones at
distinct frequencies far apart compared with the interface bandwidth add their
time-averaged forces. Statics: capillary--gravity energy of two pinned,
volume-constrained interfaces. Dynamics: inertia-free three-fluid Stokes flow
(creeping flow) on a moving axisymmetric domain. Optics: rays traced through the
modelled faces with the published indices $n=1/1.7/1$, i.e.\ as the cured lens in
air, 401 rays; Airy radii at 587.6\,nm.

\paragraph{Approximations and what is not modelled.} Non-axisymmetric ($m\ge1$)
disturbances, fluid inertia and convection, acoustic streaming (bulk and
boundary-layer driven~\cite{bach2018,jorgensen2021}), absorption heating and
thermal feedback, contact-line hysteresis, the transducers that would realize the
source modes, curing shrinkage and demoulding. The theory document notes that
streaming stresses at these pressures may be comparable to the 0.1--1\,Pa
traction demands of the lens faces. Formation runs use a linearized acoustic
load with the tangent refreshed every 10 steps.

\paragraph{Accuracy goal.} The specification is the maximum absolute height
error of each face over the clear pupil, 10\,nm; the design tool accepts a
command at 1\,nm. A small frozen residual, a convincing rendering or a converged
mesh is not a verified lens; acceptance is the chain of \cref{sec:method}.

%======================================================================
\section{Method}\label{sec:method}

\paragraph{1. Load demand.} For a face $h(r)$ pinned at $r=R$, with tension
$\sigma$ and density step $\Delta\rho$, stationarity of the graph energy at fixed
volume requires the extra normal traction
\begin{equation}
 t(r)+P=-\sigma\,\tfrac1r\bigl(r\,\phi\bigr)'+\Delta\rho\,g\,h,\qquad \phi=h'/\sqrt{1+h'^2},
 \label{eq:demand}
\end{equation}
which a prescribed face \enquote{demands} from the field up to a constant. With
$h=h_0+\tfrac12cr^2+a_4r^4+\dots$ the demand is a series
$t=\text{const}+Q_2r^2+Q_4r^4+\dots$; for a conic face its leading term vanishes
when $\Delta\rho\,g/\sigma=4kc^2$ (the conic matching rule). The checks recorded
in the study are: the finite-element meniscus under the exact demanded traction
reproduces both faces of Fig.~9b and Fig.~7a to $3\times10^{-11}$\,m; the two
discretizations of the demand agree to $2\times10^{-3}$; the series matches the
exact demand to $6\times10^{-4}$ for $r\le0.15R$; the matching rule predicts the
brine\,/\,oil choice found by brute force in S08 ($-8.3\times10^4$ vs
$-7.9\times10^4$\,m$^{-2}$, $1.0\times10^6$ vs $7.5\times10^5$\,m$^{-2}$); and
the finite-difference acoustic shape stiffness agrees with the analytic shape
Jacobian to $1.7\times10^{-5}$.

\paragraph{2. Convex capability certificate.} About the target, a tone with modal
amplitudes $c$ exerts generalized forces $F_i=c^{\mathsf H}G_ic$. For several
tones at distinct frequencies the time-averaged forces add, so with
$X_t\succeq0$ per band
\begin{equation}
 \min_{X_t\succeq0,\ \bm v}\ \bigl\lVert O(\bm F(X)-\bm d)+V\bm v\bigr\rVert_\infty
 \quad\text{s.t.}\quad \textstyle\sum_t\tr(W_tX_t)\le P_{\max},\quad
 \bm p_{t,s}^{\mathsf H}X_t\bm p_{t,s}\le p_{\max}^2,
 \label{eq:capability}
\end{equation}
is a semidefinite program; $O=HK^{-1}$ maps forces to pupil heights, $\bm d$ is
the demand and the fill-pump trims $\bm v$ act through volume columns $V$. Its
optimum lower-bounds the error of any drive in those bands. It is a relaxation in
the sense of~\cite{luo2010}; rank bounds for extreme optimal points follow
from~\cite{barvinok1995,pataki1998}. A static fold margin $\mu$ at the target
can be imposed as the linear matrix inequality
$\sym(K-A(X))\succeq\mu K$~\cite{boyd1994}. The certificate is first-order: it
ignores the acoustic shape stiffness and the path to the target.

\paragraph{3. One tone per band.} The leading eigenvectors of the optimal $X_t$
are refined locally to exactly one tone per band within both caps. If the
one-tone error exceeds the goal, the next drive option (a second band) is tried.

\paragraph{4. Start-up from the acoustics-off meniscus.} The command is ramped
as $s\,c$, $s:0\to1$, and the coupled equilibrium
$s^2F_{\mathrm{ac}}(q;c)=\nabla E(q)$ is continued from the capillary--gravity
meniscus $q_0$ at the lens's fill volumes. The smallest real part of the
eigenvalues of $K^{-1}(K-s^2A)$ along the branch is the \emph{fold margin} (1:
no acoustic softening; 0: fold). Under a slow ramp and a uniformly stable branch,
creeping-flow start-up follows this branch (theory document, Proposition
\enquote{Slow start-up follows the branch}).

\paragraph{5. Independent checks at the held command.}
(a) \emph{Fixed-command refinement}: the same physical command is re-evaluated on
wave meshes $24\times9$ to $72\times27$, element orders P3--P5, quadrature $+2$ and
36 to 72 surface elements; the first-order coupled shift $(K-A)^{-1}r$ of the
saved state is reported. No new inverse is solved.
(b) \emph{Coupled stability}: static spectrum of $K^{-1}(K-A)$, frozen-to-coupled
amplification $\lVert(K-A)^{-1}K\rVert$, and the creeping-flow spectrum on
$64\times12$ and $128\times24$ fluid meshes.
(c) \emph{Formation}: forward creeping-flow start-up from the acoustics-off
meniscus under the held command, 1\,s ramp, 2\,s total, at time steps of 25 and
12.5\,ms. Every displayed frame is a saved state; nothing is interpolated toward
the target.
The same checks also give first-order physical sensitivities at the fixed command
(\cref{tab:sensitivity}).

%======================================================================
\section{Results}\label{sec:results}

\subsection{Published lenses on the common machine}\label{sec:published}

Fig.~9b of~\cite{silva2020} has the object at infinity, $t=10$\,mm and an
aperture of 56\,mm (T/D 0.18); scaled 1/14 it is a 0.714\,mm thick lens with
f/1.42 at the image. Fig.~7a has $d_0=-150$, $d_2=70$, $t=20$\,mm and shape
factor 0.444 (T/D 0.36). \Cref{tab:published} collects the four fully checked
designs; \cref{fig:fig9b-setup,fig:fig9b-faces} show the Fig.~9b design and
\cref{fig:formation} the formation histories.

\begin{table}[htbp]\centering\small
\caption{Fully checked designs on the common machine (front\,/\,back where two
values are given). Rows marked $^\ast$ are quoted from the S09 study text; the
others were read from the run records listed in \cref{sec:provenance}. Fold
margin: minimum along the start-up branch. Decay: slowest creeping-flow rate on
the $128\times24$ mesh. Refinement: largest change of the predicted equilibrium
over all numerical trials (in brackets: excluding the P3 order).}\label{tab:published}
\begin{tabularx}{\textwidth}{@{}LRRRR@{}}
\toprule
 & Fig.~9b & Fig.~7a, unconstrained & Fig.~7a, margin $\ge0.7$ & Lens 7, two tones \\
\midrule
Acoustics-off error (\textmu m) & 0.77\,/\,2.07 & 19.4\,/\,23.1 & 19.4\,/\,23.1 & 1.10\,/\,43.8 \\
Bands (MHz) & 1.5 & 1.5 & 1.5 & 0.75 + 1.5 \\
Source power (W) & 0.17 & 0.91 & 1.12 & 0.55 \\
Sampled peak pressure (MPa) & 0.55 & 1.86 & 2.10 & 0.98 + 1.25 \\
Fold margin, start-up & 0.96 & 0.068 & 0.94 & 0.96 \\
Frozen-to-coupled amplification & 1.05 & 18.5 & 1.09 & 1.06 \\
Slowest decay (s$^{-1}$) & $-41.4$ & $-7.4$ & $-96.9$ & $-84.8$ \\
Refinement, max change (nm) & 0.18 (0.06) & 11.4 (3.8) & 2.6 (0.55) & 1.4 (0.39) \\
Formed, dt 25\,ms (nm) & 0.05\,/\,0.04 & 14.2\,/\,0.36 & 0.53\,/\,0.41 & 0.12\,/\,0.61 \\
Formed, dt 12.5\,ms (nm) & 0.05\,/\,0.04 & 10.7\,/\,0.36 & 0.53\,/\,0.41 & 0.12\,/\,0.61 \\
Spot max, formation record (\textmu m) & 0.071 & 0.71 & 0.86 & 0.125 \\
Spot RMS\,/\,max, pin-plane pupil$^\ast$ (\textmu m) & 0.004\,/\,0.071 & 0.033\,/\,0.37 & 0.027\,/\,0.45 & 0.010\,/\,0.125 \\
Airy radius at 587.6\,nm$^\ast$ (\textmu m) & 1.02 & 0.68 & 0.68 & 0.57 \\
\bottomrule
\end{tabularx}
\end{table}

\paragraph{Reading the table.} Fig.~9b needs a weak field: 0.17\,W against a
meniscus that is already within 2.1\um\ of the target, and every check moves it
by less than 0.2\,nm. Heights below about 0.2\,nm (Fig.~9b) and 0.5\,nm
(Fig.~7a) are below the model's own discretization resolution, which the
refinement rows measure. The unconstrained single tone for Fig.~7a reaches the
target but sits within 7\,\% of a fold (margin 0.068, amplification 18.5): its
refinement moves the front face by up to 11\,nm and its formation is still
decaying at the end of the run. The saved front-face error decays at a
log-slope of $-7.0$\,s$^{-1}$ between 1.5 and 2.0\,s, close to the slowest decay
rate $-7.4$\,s$^{-1}$, so the 14.2 vs 10.7\,nm difference between the two time
steps is an unsettled transient rather than two different end states. With the
fold-margin inequality at $\mu=0.7$ the refined tone has margin 0.94 and
amplification 1.09, refinement within 0.55\,nm apart from the P3 order, and both
time steps give the same formed state to the digits shown. The price is power:
1.12\,W and 2.10\,MPa, above the machine's nominal caps (both Fig.~7a syntheses
were run with limits of 10\,W and 5\,MPa).

On the coarse $64\times12$ fluid mesh every design also shows two decay rates
near $-10^{-3}$\,s$^{-1}$. S08 traced these to the 20\um\ annulus between the
clear aperture and the pinned rim, narrower than one coarse cell; they vanish on
the $128\times24$ mesh, which is the one quoted. No design has an unstable
creeping-flow mode or a negative static eigenvalue.

\begin{figure}[htbp]\centering
\includegraphics[width=\textwidth]{\figdir/figures/fig9b/s09-fig9b-setup.png}
\caption{Fig.~9b on the common machine (rendering of the saved numerical design
state, not an experiment). (a) Modelled cell, front half removed, equal scale;
(b) meridional section at equal scale with rays traced with the published
indices; (c--e) source velocity amplitude and phase of the 26-mode command on the
lower window, upper window and wall; (f) acoustic pressure at one instant,
1.5\,MHz. Source power 0.17\,W, sampled peak pressure 0.55\,MPa.}\label{fig:fig9b-setup}
\end{figure}

\begin{figure}[htbp]\centering
\includegraphics[width=\textwidth]{\figdir/figures/fig9b/s09-fig9b-faces.png}
\caption{Fig.~9b faces over the clear aperture: sag (left; Cartesian target,
acoustics-off meniscus and design overlap at this scale), height error to the
target (middle; the green band is the $\pm10$\,nm goal) and geometric spot
(right; green circle 1\um). Note the different axis scales: the acoustics-off
spot spans about $\pm6.5$\um, the design spot about 0.05\um. Cycle-mean
heights, no piston removed.}\label{fig:fig9b-faces}
\end{figure}

\begin{figure}[htbp]\centering
\includegraphics[width=\textwidth]{\genfig/formation-histories.png}
\caption{Formation in physical time for the four fully checked designs: maximum
pupil height error of each face against time, at time steps 25\,ms (solid) and
12.5\,ms (dashed), log scale. The drive ramps up over the first second (dotted
line). Generated by \file{make_figures.py} from the saved
\file{observations.json} of each formation run.}\label{fig:formation}
\end{figure}

\subsection{Sensitivities}\label{sec:sensitivities}

\Cref{tab:sensitivity} gives the coupled first-order shift of each face when one
physical parameter changes and the command is held. These are the numbers that
decide whether a physical machine could run open loop.

\begin{table}[htbp]\centering\small
\caption{Coupled first-order shift of the pupil (nm, front\,/\,back) at the held
command, from the force-convergence records ($c$: sound speed; tension: both
interfacial tensions). Lenses 15 and 71 are the two-tone
rescued designs of \cref{sec:followup}.}\label{tab:sensitivity}
\footnotesize\setlength{\tabcolsep}{4pt}
\begin{tabularx}{\textwidth}{@{}lRRRRRR@{}}
\toprule
Perturbation & Fig.~9b & Fig.~7a uncon. & Fig.~7a $\mu=0.7$ & Lens 7 & Lens 15 & Lens 71 \\
\midrule
Water $c$ $+0.1$\,\% & 28\,/\,47 & 11400\,/\,103 & 299\,/\,330 & 147\,/\,78 & 362\,/\,50 & 56\,/\,36 \\
Resin $c$ $+1$\,\% & 93\,/\,553 & 267000\,/\,3820 & 2160\,/\,3610 & 400\,/\,331 & 4470\,/\,1100 & 409\,/\,1100 \\
Frequency $+100$\,ppm & 2.0\,/\,0.3 & 1460\,/\,27 & 50\,/\,67 & 13\,/\,9.2 & 8.2\,/\,6.9 & 9.0\,/\,14 \\
Tension $+1$\,\% & 38\,/\,25 & 4270\,/\,188 & 161\,/\,113 & 27\,/\,367 & 2.3\,/\,187 & 153\,/\,83 \\
\bottomrule
\end{tabularx}
\end{table}

A 0.1\,\% change of the water sound speed (a fraction of a kelvin) moves the faces
by tens to hundreds of nanometres, and the near-fold Fig.~7a tone by 11\um. The
resin sound speed, still a placeholder, enters at the micrometre level for most
designs. These are linear predictions of the model; they say that the predicted
nanometres are only meaningful together with measured constants and closed-loop
trim against a measured wavefront.

\subsection{Atlas: 81 singlets on the common machine}\label{sec:atlas}

The family is Silva-Lora and Torres singlets with $n=1/1.7/1$; object at
infinity, $-300$ or $-150$\,mm; image at 60, 90 or 120\,mm; thickness 10, 15 or
20\,mm; shape factor $-0.5$, 0 or $+0.5$ (paper units), scaled 1/14: 81 lenses.
Thirteen do not fit the common cell (faces crossing at the rim or too close to a
window): lenses 0, 1, 2, 27, 28, 29, 54, 55, 56, 59, 63, 64 and 65.

\paragraph{First-order atlas.} \Cref{fig:atlas} shows the corrected first-order
survey (after the mesh fix of \cref{sec:mesh}): with one tone per band, 26
modes, $\le1$\,W and $\le2$\,MPa per band, a single tone at 1.5\,MHz reaches the
10\,nm goal for 58 lenses and a tone at 0.75\,MHz plus one at 1.5\,MHz for the
other 10 (lenses 6, 7, 30, 31, 33, 34, 42, 58, 60, 61); none is out of reach.

\begin{figure}[htbp]\centering
\includegraphics[width=0.92\textwidth]{\figdir/atlas-2026-10-26/figures/s09-atlas.png}
\caption{First-order atlas: least source freedom with which the refined command
reaches the 10\,nm goal for each of the 81 singlets (dark green: one tone at
1.5\,MHz; light green: two tones; grey: does not fit). Each lens is drawn at
equal scale in the common cell. This is a frozen-compliance result; the
nonlinear start-up check is \cref{fig:atlas-nonlinear}.}\label{fig:atlas}
\end{figure}

\paragraph{Nonlinear check of the whole atlas.} For every lens that fits, the
convex program and the one-tone-per-band refinement were repeated under both caps
and the coupled start-up branch was followed from the acoustics-off meniscus,
which lies 3 to 90\um\ from the target (worse face) depending on the lens (5.4 CPU-hours in
total). \Cref{tab:atlas,fig:atlas-nonlinear} summarize the end of the branch.

\begin{table}[htbp]\centering\small
\caption{End of the start-up branch, maximum pupil error over both faces, with the
command class of the first-order atlas (\file{atlas-nonlinear.jsonl}).}\label{tab:atlas}
\begin{tabular}{@{}lrl@{}}
\toprule
Start-up end error & Lenses & Detail \\
\midrule
below 1\,nm & 61 & 51 one tone, 10 two tones; median of all 68: 0.16\,nm \\
1 to 10\,nm & 3 & lenses 3, 12, 35: 2.5, 1.4, 1.1\,nm \\
10 to 20\,nm & 3 & lenses 15, 71, 78: 10.2, 12.9, 17.1\,nm \\
fold before full amplitude & 1 & lens 57: fold at $s=1$, 31\,nm \\
\bottomrule
\end{tabular}
\end{table}

\begin{figure}[htbp]\centering
\includegraphics[width=\textwidth]{\genfig/atlas-nonlinear.png}
\caption{(a) Maximum pupil error at the end of the start-up branch for each of
the 68 lenses that fit (circles: one tone; squares: two tones; cross: fold). Arrows
lead to the two-tone follow-up of the seven lenses that missed 1\,nm
(\cref{sec:followup}); for lens 57 the follow-up ends on another equilibrium.
(b) Frozen (first-order) error of the refined command against the nonlinear
start-up error. Below about 0.3\,nm the nonlinear values sit above the frozen
ones, at the level of the discretization resolution of \cref{tab:published}.
Log scales. Generated by \file{make_figures.py}.}\label{fig:atlas-nonlinear}
\end{figure}

For the seven lenses that missed 1\,nm the convex bound is 0.003--0.03\,nm (lens
57: 0.47\,nm), so the shortfall is not in the source. It is already present in
the one-tone refinement: the frozen error of the refined command is within 30\,\%
of the nonlinear one (\cref{fig:atlas-nonlinear}b), and the refined tone sits at
the 2\,MPa cap in six of the seven and at the 1\,W cap in lens 12. It is a limit
of the one-tone restriction under the caps.

\subsection{Two-tone follow-up and the lens 57 path failure}\label{sec:followup}

The seven short lenses were rerun with one tone at 0.75\,MHz plus one at 1.5\,MHz
and, separately, with one tone and a fold-margin inequality $\mu\ge0.5$
(\cref{tab:followup}).

\begin{table}[htbp]\centering\small
\caption{Follow-up of the seven lenses that missed 1\,nm (start-up end, maximum
pupil error). Two-tone column: power and the minimum fold margin along the
branch. The fold-margin run was interrupted after lenses 3, 12 and 15.}\label{tab:followup}
\begin{tabular}{@{}rrlr@{}}
\toprule
Lens & One tone, 1.5\,MHz & Two tones, 0.75 + 1.5\,MHz & One tone, $\mu\ge0.5$ \\
\midrule
3 & 2.5\,nm & 0.06\,nm (0.94\,W, margin 0.21) & 129\,nm \\
12 & 1.4\,nm & 0.03\,nm (0.69\,W, margin 0.72) & 0.23\,nm \\
15 & 10.2\,nm & 0.03\,nm (0.31\,W, margin 0.93) & 0.42\,nm \\
35 & 1.1\,nm & 0.055\,nm (0.57\,W, margin 1.00) & not run \\
71 & 12.9\,nm & 0.04\,nm (0.42\,W, margin 0.99) & not run \\
78 & 17.1\,nm & 0.03\,nm (0.53\,W, margin 0.94) & not run \\
57 & fold, 31\,nm & \textbf{other equilibrium: 25\,/\,159\um} & not run \\
\bottomrule
\end{tabular}
\end{table}

The second band rescues six of the seven, so 67 of 68 lenses form within 1\,nm
from rest. The fold-margin inequality trades accuracy for margin and is not the
cheap remedy: lens 3 gets worse (129\,nm).

\paragraph{Lens 57} (object $-150$, image 60, $t=15$, shape factor $-0.5$;
acoustics-off error 2.5\,/\,76\um) fails differently. Its two-tone command holds
the target in the frozen model (0.33\,nm; convex bound 0.019\,nm) and the branch
does reach full amplitude (minimum fold margin 0.32), but on another equilibrium,
25\,/\,159\um\ from the target, i.e.\ farther than the meniscus it started from.
Capability is not the issue; the path is. A finer uniform ramp (40 steps) is
being tested: at the time of writing \file{lens57-fine-ramp.jsonl} does not exist
and the run's log contains no result, so this test is \textbf{pending}. Ramping
the bands separately or ramping the volumes has not been tried.

\paragraph{Full checks of the rescued lenses 15 and 71.} Both were re-synthesized
(two tones) with the command \texttt{afl}~\file{capability-synthesis} and taken into the independent
checks (\cref{tab:rescued}). Note that these are separate synthesis runs: their
commands differ from those of the follow-up table (for lens 15, 0.59\,W instead
of 0.31\,W).

\begin{table}[htbp]\centering\small
\caption{Independent checks of the two-tone designs of lenses 15 and 71
(\file{rescued-2026-10-27/}).}\label{tab:rescued}
\begin{tabularx}{\textwidth}{@{}p{0.36\textwidth}LL@{}}
\toprule
 & Lens 15 & Lens 71 \\
\midrule
Lens (object, image, $t$ in paper mm; shape factor) & $\infty$, 90, 20; $-0.5$ & $-150$, 90, 20; $+0.5$ \\
Acoustics-off error (\textmu m) & 1.07\,/\,23.3 & 12.5\,/\,11.4 \\
Power (W); band peak pressures (MPa) & 0.59; 0.63 + 0.97 & 0.55; 0.59 + 0.79 \\
Start-up end (nm); fold margin & 0.058\,/\,0.207; 0.86 & 0.089\,/\,0.044; 1.00 \\
Refinement, max change (nm) & 1.0 (P3); 0.19 otherwise & 0.37 (P3); 0.12 otherwise \\
Amplification; slowest decay $128\times24$ (s$^{-1}$) & 1.19; $-77.1$ & 1.01; $-87.2$ \\
Unstable modes; negative static eigenvalues & 0; 0 & 0; 0 \\
Formation, dt 25\,ms & \textbf{pending} (last saved state $t=0.825$\,s) & \textbf{pending} (last saved state $t=0.900$\,s) \\
Formation, dt 12.5\,ms & \textbf{pending} (not started) & \textbf{pending} (not started) \\
\bottomrule
\end{tabularx}
\end{table}

The formation runs were in progress when this report was written; no
\file{validation.json} exists for either, and no result is claimed. The partial
trajectories are mid-ramp and say nothing yet about the formed state.

\subsection{Minimum hardware: channel count}\label{sec:channels}

The 112 channels discretize the continuous source. \Cref{tab:channels,fig:channels}
repeat the frozen capability and the one-tone refinement with fewer rings per
window and wall bands per layer (below 26 channels the mode count equals the
channel count).

\begin{table}[htbp]\centering\small
\caption{Channel study (first order, 1\,W, 2\,MPa; \file{channel-study.jsonl}).
Channels $=2\times$rings per window $+\,3\times$wall bands per layer. Errors are
frozen maximum pupil errors in nm.}\label{tab:channels}
\setlength{\tabcolsep}{4pt}
\begin{tabular}{@{}lrrrrrr@{}}
\toprule
 & Convex bound, & \multicolumn{5}{c}{One tone per band (refined)} \\
\cmidrule(l){3-7}
Channels (rings, wall bands) & worst of 5 & Fig.~9b & Fig.~7a & Lens 22 & Lens 45 & Lens 7 (2 tones) \\
\midrule
112 (32, 16) & 0.010 & 0.020 & 0.12 & 0.01 & 0.12 & 0.001 \\
56 (16, 8) & 0.010 & 0.005 & 0.57 & 0.02 & 0.19 & 0.001 \\
42 (12, 6) & 0.010 & 0.012 & 11.2 & 0.01 & 0.14 & 0.001 \\
28 (8, 4) & 0.011 & 0.025 & 2.8 & 0.01 & 0.04 & 0.001 \\
21 (6, 3) & 0.011 & 0.008 & 11.2 & 0.04 & 0.04 & 0.001 \\
14 (4, 2) & 0.025 & 0.020 & 12.3 & 0.34 & 0.28 & 0.007 \\
\bottomrule
\end{tabular}
\end{table}

\begin{figure}[htbp]\centering
\includegraphics[width=\textwidth]{\genfig/channel-study.png}
\caption{Channel study: (a) convex bound and (b) one-tone-per-band error against
the number of source channels, five lenses, log scales. First-order results;
none of the reduced layouts has been through the nonlinear chain.}\label{fig:channels}
\end{figure}

With 14 channels the bound is still 0.025\,nm for every lens tried and four of
the five keep a one-tone command below 0.4\,nm. Only the pressure-limited Fig.~7a
needs about 56 channels for a sub-nanometre single tone. These are first-order
numbers.

\subsection{The design tool \texttt{afl design}}\label{sec:design}

\texttt{afl design} turns a lens prescription into machine settings: it places
the lens with its front vertex 0.8\,mm above the bottom window, checks the fit,
sets the pin heights to the rim heights and the fills to the enclosed volumes,
computes the certificate, refines one tone per band (trying 1.5\,MHz, then
0.75 + 1.5\,MHz, accepting at 1\,nm) and ramps the command from rest.

\paragraph{Fig.~9b end to end.} From the prescription alone (object collimated,
image 90, $t=10$, intermediate point 137.95\,mm, scale 1/14) the tool returned
status \enquote{formed} (\file{artifacts/designs/fig9b-2026-10-27/recipe.json};
\cref{tab:recipe}). The run took 4239\,s (71\,min) on a shared 4-core machine.
The recipe reproduces the S09 nonlinear check of Fig.~9b
(\file{nonlinear-2026-10-26/fig9b-1.5mhz-26modes.json}) to all printed digits.
Two cautions: the status means that the ramp from rest ends within the goal, not
that the independent checks were run (the recipe says so); and the recipe's
command (0.1715\,W) is close to but not identical with the separately verified
S09 Fig.~9b command (0.1712\,W), so the checks of \cref{tab:published} apply to a
neighbouring command.

\begin{table}[htbp]\centering\small
\caption{Recipe of the end-to-end Fig.~9b design (SI; volumes in \textmu L $=10^{-9}$\,m$^3$).}\label{tab:recipe}
\begin{tabularx}{\textwidth}{@{}p{0.45\textwidth}L@{}}
\toprule
Pin heights above the bottom window & 1.3251\,mm and 1.5539\,mm \\
Fill volumes (lower, lens, upper) & 13.57, 6.08, 19.17\,\textmu L \\
Drive & one tone, 1.5\,MHz, 112 channels; 0.1715\,W; 0.545\,MPa \\
Convex bound; frozen error of the command & 0.0022\,nm; 0.0080\,nm \\
Start-up end error (front\,/\,back); fold margin & 0.044\,/\,0.033\,nm; 0.956 \\
Spot RMS\,/\,max; Airy radius; working f-number & 0.0034\,/\,0.054\um; 1.017\um; 1.42 \\
\bottomrule
\end{tabularx}
\end{table}

\paragraph{Lens 7.} The user guide records a \texttt{--no-startup} run for lens 7
in which one tone was tried and then two (17\,min). No recipe file of that run is
in the artifacts tree, so we give no numbers for it here; the nonlinear evidence
for lens 7 is the S09 synthesis and its full checks (\cref{tab:published}).

\subsection{What an outside listener would hear}\label{sec:sound}

For presentation only, the harmonic fields of Fig.~9b (lens A) and lens 7 (lens
B) were carried to a listener on the axis 50\,mm above the cell by a
Kirchhoff--Helmholtz integral over the revolved cell boundary, assuming a water
bath outside (the model has no exterior medium). The predicted amplitudes are
5.4\,kPa for lens A at 1.5\,MHz and 2.0 and 1.9\,kPa for lens B at 0.75 and
1.5\,MHz; doubling the listener distance multiplies them by 0.47--0.51, as
expected in the far field. The soundtrack plays each band at its frequency
divided by 2000. This is a sonification of the model, not a recording, and it
adds an assumption (the bath) that the shaping results do not use.

%======================================================================
\section{Corrections made along the way}\label{sec:corrections}

\paragraph{S08: \enquote{water cannot reach Fig.~9b} withdrawn.} S08 found that
the single-tone faithful inverse stalled at 674--676\,nm for Fig.~9b in the water
cell, at 1\,W and at 10\,W alike, and concluded that the liquids had to be tuned
(brine below, a hypothetical 5\,mN\,m$^{-1}$ oil above). That floor was a local
minimum of a non-convex inverse started from frozen least-squares fits of random
seeds. The convex program, for the same lens, cell type, 26-mode source and
frequency, gives a single tone of 0.17\,W whose start-up branch ends within
$5\times10^{-11}$\,m of both faces. The S08 multi-tone figures came from the same
kind of local fit and are equally unreliable; the tuned-liquid design remains
valid as a result, but tuning is not necessary. \emph{Why it matters:} a stalled
local optimizer is not evidence of unreachability. A reachability claim needs a
certificate, and the convex bound now serves both as that certificate and as the
initializer.

\paragraph{S09 mesh defect.}\label{sec:mesh} A first pass of the atlas used a wave
mesh that, for 20 of the 68 geometries (lenses 3, 7, 8, 15, 18, 21, 22, 25, 26,
30, 32, 38, 45, 49, 52, 58, 60, 62, 71, 77), contained zero-height elements: cell
and port grids shared layer fractions computed one rounding unit apart, so exact
de-duplication kept both rows and the wave matrix was singular (about 5\,\% linear
residual with both direct solvers; force discontinuous in the shape). Those rows
made 13 lenses look hard and 3 unreachable, and supported a claim that a second
band beats doubling the modes (14 vs 5 rescued lenses); that claim is withdrawn.
The mesh now merges near-duplicate nodes; meshes without such nodes, including
every saved S06--S09 design except the two-band atlas-38 run (marked invalid),
are bit-identical after the fix. The invalid records are kept and marked.
\emph{Why it matters:} an optimizer's small residual does not certify the forward
solve. The wave solve's linear residual (now 1--3$\times10^{-13}$ in every coupled-stability record
read for this report) belongs in every acceptance.

\paragraph{Pupil convention of the spot.} The spot values marked $^\ast$ in
\cref{tab:published} launch rays through points on the pin plane, a pupil slightly
inside the 2\,mm clear radius. Launched through the front face out to the full
clear radius, the Fig.~7a margin-0.7 maximum is 0.86\um\ (RMS 0.03\um, wavefront
PV about $\lambda/560$ at 1\,nm), slightly beyond the 0.68\um\ Airy radius at the
extreme edge ray; the Fig.~9b and lens-7 values are unchanged. The formation
records' spot maxima (0.855\um\ for the margin design, 0.705--0.712\um\ for the
unconstrained one) coincide with this full-aperture value. \emph{Why it matters:}
\enquote{diffraction-limited} holds for Fig.~9b and lens 7 under either
convention, but for Fig.~7a only for the pin-plane pupil; a spot claim must state
its pupil.

\paragraph{Radiation integral: first to second order.} The first version of the
outside-listener integral took $\partial p/\partial n$ from the gradient of the
linear interpolant on each boundary triangle, which is first-order accurate. It
was replaced by a local quadratic least-squares fit of $p$ and $\nabla p$ at each
boundary-edge midpoint (\file{src/acoustic_freeform/acoustics/radiation.py}),
tested against the exact field of a monopole inside the cell: the error is below
0.5\,\% at two listener distances and falls by more than a factor 3.5 when the
mesh is refined twofold (\file{tests/test_radiation.py}). The listener numbers of
\cref{sec:sound} were regenerated with the second-order version; the first-order
values are not in the artifacts tree, so the size of the change is not reported.
\emph{Why it matters:} it affects only the presentation soundtrack, not any
shaping result, but the test against an independent exact solution is the kind
of check the project requires.

\paragraph{Two points of bookkeeping in this report.} (i) The S09 verdict text and
the README state \enquote{61 with one tone at 1.5\,MHz, 6 more with a second
tone}. The records give 61 lenses below 1\,nm with the first command class, of
which 51 use one tone and 10 already use two tones, plus 6 rescued by a second
tone: 67 in all, 51 with one tone and 16 with two. This report uses the record
values. (ii) The Fig.~7a designs were synthesized with limits of 10\,W and
5\,MPa; the margin-0.7 design exceeds the machine's 1\,W\,/\,2\,MPa caps by 12\,\%
and 5\,\%.

%======================================================================
\section{Limits, open questions and next steps}\label{sec:limits}

\paragraph{What is established, and how.} Implemented physics: axisymmetric
lossless acoustics with ideal velocity sources, radiation stress, capillarity,
gravity, pinned rims, fixed volumes and creeping flow. Numerical verification:
fixed-command refinement of mesh, order, quadrature and surface space; coupled
static and creeping-flow stability; formation at two time steps; closed-form and
independent checks of the demand, the matching rule and the shape Jacobian. Not
done: experimental validation of any kind.

\paragraph{Limits.}
\begin{itemize}
\item \emph{Model scope.} Acoustic streaming, absorption heating and thermal
feedback, the stability of the driven state under non-axisymmetric ($m\ge1$) disturbances, inertia,
contact-line hysteresis, transducer realization of the source modes, curing
shrinkage and demoulding are not assessed. A quiescent spectrum is not the
stability of a streaming base flow.
\item \emph{Coverage of the checks.} The full chain has been run on three designs;
the start-up branch alone on all 68 atlas lenses; formation of the rescued lenses
15 and 71 is pending; the channel study is first order only.
\item \emph{Sensitivity.} Tens to hundreds of nanometres per 0.1\,\% of sound
speed; the resin sound speed and the interfacial tension are not measured.
\item \emph{Optics.} The lenses are stigmatic by construction; axial stigmatism
does not imply aplanatism or achromatism, and the spots are traced with the
published indices as the cured lens in air, not through the liquid cell.
\item \emph{Liquid optics vs.\ a product.} The results concern a driven liquid
equilibrium and its formation. Curing is a separate, unmodelled process.
\end{itemize}

\paragraph{Next steps} (in the order of the S09 verdict).
\begin{enumerate}
\item The start-up path of lens 57: the finer uniform ramp (pending), then
separate band ramps or volume ramps.
\item The full chain for a sample of the rescued lenses (15 and 71 under way) and
for reduced channel counts.
\item Acoustic streaming, heating and non-axisymmetric stability.
\item A measured NOA~61 sound speed.
\item A physical transducer realization of the source modes.
\end{enumerate}
Nothing in this report has been built or measured.

%======================================================================
\section*{Provenance of the numbers}\label{sec:provenance}
\addcontentsline{toc}{section}{Provenance of the numbers}
\begingroup\small\raggedright
Study texts: \file{studies/S09-universal-source/study.md},
\file{studies/S08-published-singlet/study.md}, \file{STATUS.md},
\file{studies/ledger.json}, \file{docs/design-a-lens.md}. Run records under
\file{artifacts/studies/S09-universal-source/} (private): synthesis
\file{synthesis-2026-10-26/<run>/validation.json}; start-up checks
\file{nonlinear-2026-10-26/*.json}; full checks in
\file{fig9b-verification-2026-10-26/}, \file{fig7a-verification-2026-10-26/},
\file{fig7a-margin0p7-verification-2026-10-26/} and
\file{atlas7-2band-verification-2026-10-26/}, each with
\file{force-convergence/}, \file{coupled-stability/}, \file{formation-dt0025/}
and \file{formation-dt00125/} (\file{validation.json}, \file{observations.json}); atlas \file{atlas-nonlinear-2026-10-27/atlas-nonlinear.jsonl},
\file{followup-two-bands.jsonl}, \file{followup-margin0p5.jsonl}; rescued lenses
\file{rescued-2026-10-27/}; channels \file{channels-2026-10-27/channel-study.jsonl};
listener \file{story-2026-10-27/story-listener.json}; design tool
\file{artifacts/designs/fig9b-2026-10-27/recipe.json}. Values quoted only from
the study text (marked $^\ast$ in \cref{tab:published}, the theory checks of
\cref{sec:method}, and the full-aperture spot RMS and wavefront PV) were not found
in a separate run record. Configurations: \file{configs/S09-universal-source/}.
Generated plots: \file{docs/reports/s09-results/make_figures.py}.
\par\endgroup

\printbibliography[heading=bibintoc,title={References}]
\end{document}
